\documentclass[12pt]{iopart}
\usepackage{graphicx}
\usepackage{hyperref}
\usepackage{cite}
\usepackage{booktabs}
\usepackage{qrcode}
\usepackage{amssymb}

% --- iopart definiert einige Standardkommandos nicht ---
\providecommand{\mathbb}[1]{\mathbf{#1}}

\begin{document}
	
	\title{Resonanzfeldtheorie: Axiomatik, Grundformel und
		empirische Validierung}
	
	\author{Dominic-Ren\'e Schu}
	\address{Unabh\"angiger Forscher, Deutschland\\
		\href{https://github.com/DominicReneSchu/public}%
		{https://github.com/DominicReneSchu/public}\\
		ORCID: 0009-0004-9769-9061\\
		Email: dominic.rene.schu@gmail.com}
	
	\begin{abstract}
		Die Resonanzfeldtheorie (RFT) wird als axiomatisches Rahmenwerk
		vorgestellt, das Schwingungskopplung als universelles
		Organisationsprinzip physikalischer und systemischer Ph\"anomene
		beschreibt. Aus sieben Axiomen folgt eine parametrische Erweiterung
		der Planck-Relation um einen phasenabh\"angigen Kopplungsoperator
		$\varepsilon(\Delta\varphi) \in [0,1]$:
		$E = \pi \cdot \varepsilon(\Delta\varphi) \cdot \hbar \cdot f$.
		Die Theorie wird an zwei unabh\"angigen Datens\"atzen
		empirisch validiert: (1)~Monte-Carlo-Simulation auf
		CMS-Open-Data-Dielectron-Daten, die vier bekannte
		Teilchenresonanzen (J/$\psi$, $\Upsilon$(1S), $\Upsilon$(2S),
		Z-Boson) mit empirischem p-Wert~$= 0$ \"uber $10\,000$
		Simulationen detektiert; (2)~ResoTrade, ein
		resonanzfeldtheoretisches Trading-System, das \"uber 24~Monate
		in vier Marktregimen $+26{,}1\%$ gegen\"uber HODL erzielt.
		Alle Herleitungen, Quellcodes und Daten sind offen verf\"ugbar
		unter \url{https://github.com/DominicReneSchu/public}.
	\end{abstract}
	
	\noindent\textbf{Schl\"usselbegriffe:} Resonanz, Feldtheorie,
	Kopplungsoperator, Axiomatik, Monte-Carlo-Simulation,
	empirische Validierung, Open Science
	
	\medskip
	
	\noindent\textbf{Daten- und Codeverf\"ugbarkeit:}
	\url{https://github.com/DominicReneSchu/public}
	
	
	%% ============================================================
	\section{Einleitung}
	%% ============================================================
	
	Die moderne Physik verf\"ugt \"uber Formalismen von enormer
	Vorhersagekraft, doch ihre begrifflichen Grundlagen bleiben
	fragmentiert. Quantenmechanik, Relativit\"atstheorie und
	Thermodynamik operieren in getrennten theoretischen R\"aumen.
	Programme der Quantengravitation erh\"ohen die mathematische
	Komplexit\"at, ohne empirische \"Uberpr\"ufbarkeit zu schaffen.
	
	Dieses Manuskript f\"uhrt die \textit{Resonanzfeldtheorie}
	(RFT) ein -- ein axiomatisches Rahmenwerk, das
	Schwingungskopplung als prim\"ares Organisationsprinzip
	physikalischer Systeme postuliert. Die RFT formuliert sieben
	Axiome, aus denen eine parametrische Erweiterung der
	Quantenmechanik, eine vereinheitlichte
	Differentialgleichungsstruktur und pr\"ufbare Vorhersagen
	folgen.
	
	Im Unterschied zu bestehenden Ans\"atzen zielt die RFT nicht
	auf eine Vereinheitlichung der fundamentalen Kr\"afte, sondern
	auf einen konzeptuellen Rahmen, der Resonanzph\"anomene
	\"uber Skalen und Dom\"anen hinweg einheitlich beschreibt --
	von Teilchenphysik bis Finanzm\"arkte.
	
	Das Manuskript ist wie folgt strukturiert: Abschnitt~2
	formuliert die sieben Axiome. Abschnitt~3 entwickelt die
	mathematische Struktur. Abschnitt~4 pr\"asentiert die
	empirische Validierung. Abschnitt~5 diskutiert den Vergleich
	mit etablierten Theorien. Abschnitt~6 benennt offene Fragen.
	Abschnitt~7 fasst zusammen.
	
	
	%% ============================================================
	\section{Axiomatische Grundlagen}
	%% ============================================================
	
	Die RFT basiert auf sieben hierarchisch aufgebauten Axiomen.
	Die Axiome sind nicht unabh\"angig, sondern konstitutiv
	verschr\"ankt: A1--A3 definieren die Schwingungsstruktur,
	A4 die Energetik, A5--A6 die Dynamik, A7 die Stabilit\"at.
	
	\subsection{Axiom 1: Universelle Schwingung}
	
	Jede physikalische Entit\"at ist durch eine
	Schwingungsfunktion beschreibbar:
	\begin{equation}
		\psi(x, t) = A \cos(kx - \omega t + \phi)
		\label{eq:oscillation}
	\end{equation}
	
	\textit{Beispiel:} Quantisierte Elektronenbahnen, Pendel mit
	Eigenfrequenz, zirkadiane Rhythmen, Marktzyklen.
	
	\subsection{Axiom 2: Superposition und Interferenz}
	
	Schwingungen \"uberlagern sich linear:
	\begin{equation}
		\Psi(x, t) = \sum_i \psi_i(x, t)
		\label{eq:superposition}
	\end{equation}
	
	\textit{Beispiel:} Doppelspalt-Interferenz,
	Schwebungsph\"anomene, AC/DC-Zerlegung von Preisfeldern.
	
	\subsection{Axiom 3: Resonanzbedingung}
	
	Kopplung zwischen zwei Systemen tritt auf, wenn ihre
	Frequenzen in rationalem Verh\"altnis stehen:
	\begin{equation}
		\frac{f_1}{f_2} = \frac{n}{m}, \quad
		n, m \in \mathbb{Z}^+
		\label{eq:resonance_condition}
	\end{equation}
	
	\textit{Beispiel:} Molekulare Orbitalresonanz,
	Tacoma-Narrows-Br\"ucke. Negativtest: Altcoins ohne
	Eigenfrequenz erzeugen keine Resonanz (200\,000 Episoden).
	
	\subsection{Axiom 4: Kopplungsenergie}
	
	Die durch Resonanz \"ubertragene Energie ist:
	\begin{equation}
		E = \pi \cdot \varepsilon(\Delta\varphi) \cdot \hbar \cdot f
		\label{eq:grundformel}
	\end{equation}
	mit dem phasenabh\"angigen Kopplungsoperator
	$\varepsilon(\Delta\varphi) \in [0,1]$, wobei
	$\Delta\varphi$ die Phasendifferenz zwischen gekoppelten
	Systemen bezeichnet.
	
	\textbf{Dimensionsanalyse:}
	$[\,] \cdot [\,] \cdot [\mathrm{J{\cdot}s}] \cdot
	[\mathrm{s}^{-1}] = [\mathrm{J}]$ $\checkmark$
	
	\textbf{Spezialfall Planck:}
	F\"ur $\varepsilon = 1/(2\pi)$ folgt
	$E = \frac{1}{2}\hbar f$, konsistent mit der
	Grundzustandsenergie des harmonischen Oszillators.
	
	\textbf{Grenzwerte:}
	\begin{itemize}
		\item $\varepsilon = 0$: keine Kopplung
		(destruktive Interferenz)
		\item $\varepsilon = 1$: maximale Kopplung
		(vollst\"andige Phasenkoh\"arenz)
	\end{itemize}
	
	\subsection{Axiom 5: Energierichtung}
	
	Energie ist eine vektorielle Gr\"o\ss{}e mit Betrag und
	Richtung im Resonanzfeld:
	\begin{equation}
		\vec{E}_{\mathrm{dir}} =
		E_{\mathrm{kurz}} - E_{\mathrm{lang}}
		\label{eq:energy_direction}
	\end{equation}
	
	Der Energierichtungsvektor bestimmt die Dynamik des
	Energieflusses. Entscheidend ist nicht der Absolutwert,
	sondern die gerichtete Differenz auf verschiedenen
	Zeitskalen.
	
	\textit{Beispiel:} In ResoTrade bestimmt
	\texttt{energy\_dir = e\_short - e\_long} die
	Handelsrichtung -- Phasenerkennung statt Preisprognose.
	
	\subsection{Axiom 6: Informationsfluss durch
		Resonanzkopplung}
	
	Information propagiert ausschlie\ss{}lich entlang
	koh\"arenter Resonanzpfade. Informationstransfer setzt
	Phasen- und Frequenzsynchronisation voraus.
	\begin{equation}
		\mathrm{Informationsfluss} \iff
		\mathrm{Phasen\textrm{-} und\;Frequenzsynchronisation}
		\label{eq:information}
	\end{equation}
	
	\textit{Beispiel:} Taktsignale in digitalen Systemen,
	synaptische Kopplung, Resonanz-Gate in ResoTrade.
	
	\subsection{Axiom 7: Invarianz}
	
	Resonanzstrukturen bleiben unter Variation der
	Beobachterperspektive und Systemparameter erhalten:
	\begin{equation}
		\forall\; T : \mathcal{R} \to \mathcal{R}', \quad
		\mathrm{Struktur}(\mathcal{R})
		= \mathrm{Struktur}(\mathcal{R}')
		\label{eq:invariance}
	\end{equation}
	
	\textit{Beispiel:} Monte-Carlo-Ergebnisse stabil \"uber
	$10\,000$ Durchl\"aufe. ResoTrade konsistent \"uber
	vier Marktregime.
	
	
	%% ============================================================
	\section{Mathematische Formulierung}
	%% ============================================================
	
	\subsection{Grundformel und Kopplungsoperator}
	
	Die zentrale Gleichung~(\ref{eq:grundformel}) erweitert die
	Planck-Relation $E = \hbar\omega$ um einen
	Kopplungsterm. Der Faktor $\pi$ reflektiert die
	zyklische Geometrie der Resonanzkopplung. Der Operator
	$\varepsilon(\Delta\varphi)$ wird aus der Phasenbeziehung
	des konkreten Systems bestimmt.
	
	F\"ur ein System mit Phasendifferenz $\Delta\varphi$
	zwischen Anregung und Antwort ist das einfachste Modell:
	\begin{equation}
		\varepsilon(\Delta\varphi) =
		\frac{1}{2}(1 + \cos\Delta\varphi)
		\label{eq:epsilon_cos}
	\end{equation}
	Dies liefert $\varepsilon = 1$ bei $\Delta\varphi = 0$
	(Resonanz) und $\varepsilon = 0$ bei
	$\Delta\varphi = \pi$ (Antiresonanz).
	
	\subsection{Resonanz-Differentialgleichung (rDGL)}
	
	Die RFT vereinheitlicht klassische Differentialgleichungen
	als Spezialf\"alle einer allgemeinen
	Resonanzkopplungsrelation:
	\begin{equation}
		\mathcal{R}(x, \dot{x}, \ddot{x}, t, \Phi) = 0
		\label{eq:rdgl}
	\end{equation}
	wobei $\Phi$ den Kopplungsterm beschreibt. Spezialf\"alle:
	
	\medskip
	\begin{center}
		\begin{tabular}{lll}
			\hline
			\textbf{System} & \textbf{Kopplung $\Phi$}
			& \textbf{Gleichung} \\
			\hline
			Harmon.\ Oszillator & $-\omega_0^2 x$
			& $\ddot{x} + \omega_0^2 x = 0$ \\
			Ged\"ampft & $-2\gamma\dot{x} - \omega_0^2 x$
			& $\ddot{x} + 2\gamma\dot{x} + \omega_0^2 x = 0$ \\
			Van-der-Pol & $-\mu(1{-}x^2)\dot{x} - x$
			& Selbsterregte Schwingung \\
			Stochastisch & $\Phi_{\det} + \xi(t)$
			& Langevin-Typ \\
			Netzwerk & $\sum_j K_{ij}\sin(\theta_j{-}\theta_i)$
			& Kuramoto-Modell \\
			\hline
		\end{tabular}
	\end{center}
	
	\subsection{Resonanzrelation und Feldklassen}
	
	Die bin\"are Resonanzrelation
	$R \subseteq \mathcal{G} \times \mathcal{G}$ ist eine
	\"Aquivalenzrelation (reflexiv, symmetrisch, transitiv),
	die eine Resonanzgruppe~$\mathcal{G}$ in disjunkte
	Resonanzklassen unterteilt.
	
	\subsection{Zustandsentwicklung}
	
	Die zeitliche Entwicklung des Resonanzfeldes folgt:
	\begin{equation}
		i\hbar \frac{\partial}{\partial t} |\psi(t)\rangle
		= \hat{H}_{\mathrm{res}} |\psi(t)\rangle
		\label{eq:schroedinger_res}
	\end{equation}
	mit dem Resonanz-Hamiltonoperator:
	\begin{equation}
		\hat{H}_{\mathrm{res}} = \hat{H}_0
		+ \varepsilon(\Delta\varphi)\,
		\hat{V}_{\mathrm{Kopplung}}
		\label{eq:h_res}
	\end{equation}
	wobei $\hat{H}_0$ der freie Hamiltonoperator und
	$\hat{V}_{\mathrm{Kopplung}}$ das Kopplungspotential
	ist. F\"ur $\varepsilon = 0$ entkoppelt das System;
	f\"ur $\varepsilon = 1$ ist die Kopplung maximal.
	
	
	%% ============================================================
	\section{Empirische Validierung}
	\label{sec:empirics}
	%% ============================================================
	
	Die RFT wird an zwei unabh\"angigen Dom\"anen gepr\"uft:
	Teilchenphysik (Abschnitt~\ref{sec:monte_carlo}) und
	Finanzm\"arkte (Abschnitt~\ref{sec:resotrade}).
	
	\subsection{Monte-Carlo-Simulation auf CMS-Dielectron-Daten}
	\label{sec:monte_carlo}
	
	\subsubsection{Datensatz und Methodik.}
	CMS-Open-Data-Dielectron-Datensatz~\cite{cms_opendata}:
	$99\,915$ Ereignisse, invariante Masse~$M$ im Bereich
	$2{,}00$--$110{,}00$~GeV.
	
	Die Hintergrundverteilung wird mittels Kerndichtesch\"atzung
	(KDE, Bandbreite $0{,}5$~GeV) aus den Daten extrahiert,
	unter schmalem Ausschluss der Signalbereiche
	(J/$\psi$: $\pm 0{,}15$~GeV; $\Upsilon$: $\pm 0{,}20$~GeV;
	Z: $\pm 4{,}0$~GeV). Erwartete Trefferraten werden
	$\Delta$-abh\"angig aus dem KDE-Modell vorberechnet.
	
	F\"ur jedes der $10\,000$ Pseudo-Experimente wird die
	vollst\"andige Resonanzanalyse durchgef\"uhrt: Binomialtest
	mit Bonferroni-Korrektur \"uber 68~Fensterbreiten
	($\Delta = 0{,}02$--$10{,}0$~GeV).
	Bootstrap-Konfidenzintervalle ($5\,000$~Wiederholungen)
	quantifizieren die statistische Stabilit\"at.
	
	\subsubsection{Ergebnisse.}
	
	\begin{table}[ht]
		\centering
		\caption{Monte-Carlo-Ergebnisse ($10\,000$ Simulationen,
			26.02.2026). Empirischer p-Wert~$= 0$: In keinem
			Pseudo-Experiment wurde ein vergleichbarer \"Uberschuss
			durch reinen Hintergrund erzeugt.}
		\label{tab:mc_results}
		\begin{tabular}{llrrrl}
			\hline
			$M_0$ (GeV) & Resonanz & $\Delta_{\mathrm{opt}}$ (GeV)
			& Hits & $p_{\mathrm{corr}}$ & emp.\ $p$ \\
			\hline
			$3{,}10$ & J/$\psi$ & $0{,}440$ & $3256$
			& $1{,}0 \times 10^{-121}$ & $0$ \\
			$9{,}46$ & $\Upsilon$(1S) & $0{,}780$ & $4186$
			& $3{,}0 \times 10^{-201}$ & $0$ \\
			$10{,}02$ & $\Upsilon$(2S) & $0{,}840$ & $4625$
			& $5{,}2 \times 10^{-188}$ & $0$ \\
			$91{,}20$ & Z-Boson & $0{,}040$ & $52$
			& $< 10^{-300}$ & $0$ \\
			\hline
		\end{tabular}
	\end{table}
	
	Alle vier Resonanzen werden mit extremer Signifikanz
	detektiert (Tabelle~\ref{tab:mc_results}). Dies
	best\"atigt Axiom~3 (Resonanzbedingung) und Axiom~7
	(Invarianz \"uber Simulationsdurchl\"aufe).
	
	\begin{figure}[ht]
		\centering
		\includegraphics[width=0.7\textwidth]%
		{figures/hist_mc_vs_real_hits.png}
		\caption{Trefferzahlen im optimalen Fenster: blaue
			Histogramme zeigen die MC-Verteilung, rote Linien
			die Realdaten.}
		\label{fig:mc_hist}
	\end{figure}
	
	\begin{figure}[ht]
		\centering
		\includegraphics[width=0.7\textwidth]%
		{figures/pvalue_curves.png}
		\caption{Korrigierte p-Werte als Funktion der
			Fensterbreite~$\Delta$. Blaues Band:
			68\%-Intervall der MC-Simulationen.
			Rote Kurve: Realdaten.}
		\label{fig:mc_pvalue}
	\end{figure}
	
	\begin{figure}[ht]
		\centering
		\includegraphics[width=0.7\textwidth]%
		{figures/heatmaps_hits.png}
		\caption{Heatmaps der Trefferzahlen f\"ur alle
			$(M_0, \Delta)$-Kombinationen. Links: Realdaten.
			Rechts: MC-Mittelwert.}
		\label{fig:mc_heatmap}
	\end{figure}
	
	\subsection{ResoTrade: Resonanzfeldtheorie im BTC-Markt}
	\label{sec:resotrade}
	
	ResoTrade~\cite{rftrepo} wendet die Axiome der RFT auf den
	BTC-Markt an. Das System zerlegt den Preis in
	DC-Komponente (langfristiger Trend, $\mathrm{MA}_{168h}$)
	und AC-Komponente (handelbare Schwingung) gem\"a\ss{}
	Axiom~1. Handelsentscheidungen basieren auf dem
	Energierichtungsvektor (Axiom~5) und einem
	Resonanz-Gate (Axiom~6), das Trades nur bei
	Phasenkoh\"arenz zul\"asst.
	
	\begin{table}[ht]
		\centering
		\caption{ResoTrade V11.1 Performance \"uber 24~Monate.
			Seit Februar 2026 live \"uber Kraken.}
		\label{tab:resotrade}
		\begin{tabular}{llrl}
			\hline
			Marktphase & Zeitraum & Trades & vs.\ HODL \\
			\hline
			Sideways (60k--70k) & M\"ar--Sep 2024
			& 437 & $+33{,}3\%$ \\
			Bullrun (60k--110k) & Sep 2024--M\"ar 2025
			& 429 & $+19{,}8\%$ \\
			Korrektur & M\"ar--Sep 2025
			& 247 & $+4{,}3\%$ \\
			Crash (110k$\to$63k) & Sep 2025--Feb 2026
			& 279 & $+46{,}8\%$ \\
			\hline
			\textbf{Durchschnitt} & \textbf{24 Monate}
			& \textbf{1392} & $\mathbf{+26{,}1\%}$ \\
			\hline
		\end{tabular}
	\end{table}
	
	Kein klassischer Indikator (RSI, MACD, Bollinger) erreicht
	auf demselben Datensatz eine Korrelation \"uber~$0{,}05$.
	Die Altcoin-Analyse best\"atigt Axiom~3 negativ: Systeme
	ohne Eigenfrequenz erzeugen keine Resonanz -- empirisch
	nachgewiesen \"uber $200\,000$~Episoden.
	
	\textbf{Axiom-Zuordnung:}
	\begin{itemize}
		\item A1: AC/DC-Zerlegung (Preis $=$ DC $+$ AC)
		\item A4: Balance-Regler (Cash $\leftrightarrow$ BTC)
		\item A5: Energierichtung
		(\texttt{energy\_dir = e\_short - e\_long})
		\item A6: Resonanz-Gate (Trades nur in Phase)
	\end{itemize}
	
	
	%% ============================================================
	\section{Vergleich mit etablierten Theorien}
	%% ============================================================
	
	Die RFT ist eine parametrische Erweiterung der
	Planck-Relation, kein Ersatz der Quantenmechanik.
	
	\medskip
	\begin{center}
		\begin{tabular}{lll}
			\hline
			\textbf{Aspekt} & \textbf{Standard-QM}
			& \textbf{RFT} \\
			\hline
			Energieformel & $E = \hbar\omega$
			& $E = \pi\varepsilon\hbar f$ \\
			Kopplung & Extern (St\"orungstheorie)
			& Intrinsisch ($\varepsilon$) \\
			Phasenbezug & Implizit
			& Explizit ($\Delta\varphi$) \\
			Beobachter & Extern (Messpostulat)
			& Resonator (A7) \\
			Vorhersagetyp & Wahrscheinlichkeiten
			& Richtung + Phase \\
			\hline
		\end{tabular}
	\end{center}
	
	\medskip
	F\"ur $\varepsilon = 1/(2\pi)$ reduziert sich die
	RFT-Grundformel auf den Planck-Spezialfall.
	
	
	%% ============================================================
	\section{Offene Fragen und Ausblick}
	%% ============================================================
	
	\begin{enumerate}
		\item \textbf{Unabh\"angige Bestimmung von
			$\varepsilon$:}
		Gleichung~(\ref{eq:epsilon_cos}) ist ein Modellansatz.
		Eine ab-initio-Herleitung aus der Systemgeometrie
		steht aus.
		
		\item \textbf{Neue Vorhersagen:}
		Die RFT muss eine bisher unbekannte Resonanz
		vorhersagen, die experimentell best\"atigt wird.
		
		\item \textbf{Live-Validierung:}
		ResoTrade l\"auft seit Februar 2026 live.
		6--12~Monate verifizierbare Live-Performance
		bilden den n\"achsten Pr\"ufstein.
		
		\item \textbf{Vergleich mit ML-Systemen:}
		Direkter Benchmark gegen state-of-the-art
		ML-Trading auf identischen Daten.
		
		\item \textbf{Weitere Dom\"anen:}
		\"Ubertragbarkeit auf Akustik, Biologie,
		Netzwerktheorie, gekoppelte Oszillatoren.
		
		\item \textbf{$\hat{H}_{\mathrm{res}}$ konkretisieren:}
		Der Resonanz-Hamiltonoperator~(\ref{eq:h_res}) muss
		f\"ur spezifische Systeme explizit konstruiert
		werden.
	\end{enumerate}
	
	
	%% ============================================================
	\section{Schlussfolgerung}
	%% ============================================================
	
	Die Resonanzfeldtheorie formuliert ein axiomatisches
	Rahmenwerk auf Basis von sieben Axiomen, das
	Schwingungskopplung als universelles Organisationsprinzip
	beschreibt. Die zentrale Grundformel
	$E = \pi \cdot \varepsilon(\Delta\varphi) \cdot
	\hbar \cdot f$ ist dimensionskonsistent, enth\"alt die
	Planck-Relation als Spezialfall und wird an zwei
	unabh\"angigen Dom\"anen empirisch validiert.
	
	Die Monte-Carlo-Simulation detektiert vier bekannte
	Teilchenresonanzen mit empirischem p-Wert~$= 0$.
	ResoTrade \"ubertrifft HODL in vier Marktregimen \"uber
	24~Monate. Beide Ergebnisse sind konsistent mit den
	Axiomen der RFT -- Teilchenphysik best\"atigt Axiom~3
	und~7, ResoTrade best\"atigt Axiom~1, 4, 5 und~6.
	
	Die Theorie ist in ihrer aktuellen Form ein
	konzeptionelles Rahmenwerk mit erstem nicht-trivialem
	empirischem Erfolg. F\"ur wissenschaftliche Anerkennung
	als physikalische Theorie fehlen: eine quantitative
	Vorhersage, die vom Standardmodell abweicht; eine
	unabh\"angige Methode zur Bestimmung von~$\varepsilon$
	ohne Kenntnis des Ergebnisses; und mehrj\"ahrige
	Live-Validierung von ResoTrade.
	
	Alle Herleitungen, Quellcodes und Daten sind
	vollst\"andig offen verf\"ugbar.
	
	
	\section*{Code- und Datenverf\"ugbarkeit}
	
	S\"amtliche Quellcodes, Datens\"atze und
	Simulationsskripte stehen unter
	\url{https://github.com/DominicReneSchu/public}
	zur Verf\"ugung.
	
	\medskip
	\noindent
	\qrcode[height=1.5cm]{https://github.com/DominicReneSchu/public}
	\hfill
	\begin{minipage}[b]{0.7\linewidth}
		\small QR-Code f\"ur direkten Zugang zum Repository.
	\end{minipage}
	
	
	\begin{thebibliography}{99}
		
		\bibitem{planck1901}
		Planck M 1901 \textit{Annalen der Physik} \textbf{4} 553
		
		\bibitem{einstein1905}
		Einstein A 1905 \textit{Annalen der Physik} \textbf{17} 891
		
		\bibitem{bohr1913}
		Bohr N 1913 \textit{Philosophical Magazine} \textbf{26} 1
		
		\bibitem{kuramoto1984}
		Kuramoto Y 1984 \textit{Chemical Oscillations, Waves, and
			Turbulence} (Berlin: Springer)
		
		\bibitem{cms_opendata}
		CMS Collaboration 2016 \textit{CMS Open Data Portal},
		\url{https://opendata.cern.ch}
		
		\bibitem{rftrepo}
		Schu D-R 2025/2026 \textit{Resonanzfeldtheorie: Axiomatik,
			Quellcode und empirische Validierung},
		\url{https://github.com/DominicReneSchu/public}
		
	\end{thebibliography}
	
\end{document}